\begin{scanpage}{0140}{122}
Let us consider the quantity $\widetilde g(1)$ in the Expression (21.11) to be a function of the angle $\phi$ [see Equation (20.05)], that is, let us replace $\widetilde g(1)$ by the function $g(1)$. In this case the equation
\[
\begin{aligned}
H_z(1-2)={}&\\
={}&H_{0z}\frac2\pi g(1)\times\int_\infty^{2\sqrt{ka}\cos(\pi/4-\varphi/2)}e^{iq^2}\,dq\,\frac{e^{ikr}}{\sqrt{2kr}}e^{ika(\sin\alpha-\sin\varphi)}\\
&+H_{0z}g(1)\frac{e^{i(kr+\pi/4)}}{\sqrt{2\pi kr}}e^{ika(\sin\alpha-\sin\varphi)}.
\end{aligned}\tag{21.12}
\]
will give qualitatively correct results not only when $\varphi\approx-\pi/2$, but also with all other values of $\phi$. Actually, the Fresnel integral is close to zero if $\sqrt{ka}\cos(\pi/4-\varphi/2)\gg1$, and in Equation (21.12) only the second term remains, as must be the case. Therefore, Equation (21.12) may be investigated as an interpolation equation, and it may be applied with any values of $\varphi$ ($|\varphi|<\pi/2$). It is easy to establish that now the fringing field does not undergo a discontinuity with the passage through plane $x=0$, since Expression (21.12) becomes zero when $\varphi=\pm\pi/2$.

It is interesting to note that Equation (21.12) automatically follows from Equation (21.08) if in the latter equation one replaces $\widetilde g(1)$ by $g(1)$. Essentially, this substitution is equivalent to the assumption that the moments of the filaments, the fields of which are used for approximating the primary edge waves, depend on the radiation direction (that is, on the azimuth $\phi$ of the observation point)
\[
m_{1z}^-=-m_{1z}^+=\frac1{4\pi k^2}H_{0z}g(1)e^{ika\sin\alpha}.\tag{21.13}
\]
Such a determination of the moments of the auxiliary linear sources is used, for example, in the work of Millar [47].

Precisely in the same way that Equations (21.06) and (21.07) were obtained, we find (when $x>0$)
\end{scanpage}
